\section{Implementation Details}

We use MINTO 3.1 and CPLEX 9.1 to implement our branch-and-price
algorithm. MINTO (\citet{Savelsbergh2}) is a mixed integer solver 
that uses a branch-and-bound algorithm with an LP solver. 
MINTO has built-in applications such as preprocessing, constraint 
generation and primal heuristics, but it also allows a user to write 
his/her own application code that is specific to a problem. In 
implementing our branch-and-price algorithm, we have written routines 
to generate an initial set of columns, to solve the pricing problem 
and to select a branching variable. 

\begin{comment}
{\bf Initial Feasible Solution} Before beginning the branch-and-price 
algorithm, we need to define a restricted master problem with 
an initial set of columns that contains a feasible solution. For each 
facility, we create all feasible paths (considering the vehicle capacity 
constraint and the time limit constraint) that have one customer, two 
customers and three customers. Since the cost matrix is symmetric, only
$N \choose 2$ paths with two customers are considered. For paths that include 
three customers, we select the minimum cost path among the feasible orderings 
of the customers.
\end{comment}

{\bf Pricing Problem} In step 1 of the 2-step branch and price  algorithm, we solve the 
pricing problem heuristically in which we use modified label correcting
algorithm for modified ESPPRC with label limit. We use label limit 400
at each node. 
In step 2, at each node, first we solve the pricing problem 
heuristically using ESPPRC algorithm with label limit 400. Then, if
the algorithm cannot find a column with negative reduced cost, the 
algorithm is run for label limit 800. If the heuristic ESPPRC algorithm
with label limit 800 cannot find any columns, we solve exact ESPPRC.
At each step of the column generation, we select the  
40 columns with the most negative reduced costs to add to the RMP. 

{\bf Primal Feasible Solution} In step 1 of the 2-step branch and price algorithm,
at the root node of the branch-and-price tree, we use CPLEX with a time limit of
1 hour to solve an integer program defined over the set of generated columns  
to get an incumbent solution.  
In other nodes of the branch-and-price tree, we call the default primal heuristic 
in MINTO to find a primal feasible solution. In step 2, we use the solution
obtained in step 1 as an intial upper bound and use  default primal heuristic 
in MINTO during the algorithm.

{\bf Branching Variable} At each node of the branch-and-price tree with 
a fractional solution, we first branch on any fractional location variables. 
If there is more than one fractional $T_j$ variable, we select 
the one that has a value closest to $0.5$. If all of the $T_j$ variables are 
integer, we next branch on a fractional $nV_j$ variable, which equals the number 
of vehicles in use at facility $j$. Of all the fractional $nV_j$ variables, 
we choose the one that is closest to $0.5$. In a fractional solution, if all 
the location variables and the total number of vehicle variables are integer, then we 
select the facility $j$ and the customer $i$ such that 
$\sum_{p \in P_j}a_{ipj}Z_{jp} $  is closest to 0.5. Similarly, for
branching rule 4, we choose a pair customer of nodes $(n_1,n_2)$ such that   
$\sum_{j \in M}\sum_{\{p \in P_j | (n_1\to n_2) \in p\}} Z_{jp}$
is closest to 0.5. 

{\bf Artificial variables} At every node of the branch-and-price tree, we need to start 
the column generation algorithm with a feasible solution or we need to know that the 
node is infeasible to be able to fathom it. Since the RMP does not contain all of the columns,
it is not necessarily true that the node is infeasible if the current LP at the node is 
infeasible. At the root node, we start with a feasible LP solution, but the RMP might become 
infeasible with respect to the current set of columns as the branching rules are added. 
Therefore, if the current LP solution at a node is infeasible, we need to verify that it 
truly is infeasible or we need to find a feasible solution from which to start
the column generation algorithm.    

At a node where the branching rules apply only to the facility location variables, 
it is easy to detect the infeasibility with respect to the number of facilities. By 
considering fixed and unfixed location variables,  we can simply check whether the 
lower bound on the number of selected facilities (constraint (\ref{lb-open-fac})) 
is satisfied or not.

Detecting the infeasibility  of the model with respect to the number of vehicles at 
each facility is more difficult.  One reason is that the number of pairing variables, 
the summation of which determines the number of vehicles in each facility, is large.
Another reason is that to decide whether there is a feasible solution with given upper bounds and
lower bounds for the number of vehicles in facilities, one must check the set partitioning constraints
(\ref{spp_con1}) as well as the facility capacity constraints (\ref{spp_con2}). 
In this case, we need to introduce artificial 
variables. We can add the artificial variables at the root node or at any node in which the
RMP is infeasible. Currently, we introduce the artificial variables at the root node. With the artificial 
variables, we are able to 
find an initial feasible solution at any node for any set of restrictions on the 
facility variables, the total number of vehicle variables and the pairing variables.    

We create an imaginary facility with $N$ pairings, one for each customer. Each pairing (column) 
visits only one customer. We define the following variables and parameters: 
%\singlespace
\begin{align*}
T_A & = \text {artificial facility} \nonumber \\
Z_{A0},Z_{A1},..,Z_{A(N-1)} & = \text {artificial columns} \nonumber \\
C_A & = \text {cost of an artificial column} \nonumber \\
FC_A & = \text {fixed cost of the artificial facility }\nonumber \\
l_j^u & = \text {lower bound for the number of vehicles for facility j at node u} \nonumber \\
u_j^u & = \text {upper bound for the number of vehicles for facility j at node u} \nonumber  
\end {align*}  

We add these artificial variables to our RMP: 
\begin{optprog}
Minimize & \objective{\sum_{j \in M} FC T_{j} + \sum_{j \in M} \sum_{p\in P_{j}} C_{jp} Z_{jp} + FC_A T_{A}+ C_A \sum_{i \in N}Z_{Ai}} \\
subject to & \sum_{j \in M} \sum_{p \in P_{j}} a_{ipj} Z_{jp} + Z_{Ai} & = & 1 & \forall i \in N  \label {spc}\\
           & \sum_{p \in P_j} \sum_{i \in N} a_{ipj} d_i Z_{jp} & \leq & CapF_{j} T_{j} & \forall j \in M  \\
           & \sum_{j \in M} T_{j} & \geq & nREQ \\
           & \sum_{p \in P_{j}}a_{ipj}Z_{jp}\ & \leq & T_{j} & \forall j \in M, \forall i \in N \\
           & Z_{Ai} & \leq & T_{A}  & \forall i \in N \\
           & \sum_{p \in P_j}Z_{jp} & \leq & l_j^u  & \forall j \in M \label{low-bd}\\
           & \sum_{p \in P_j}Z_{jp} & \geq & u_j^u  & \forall j \in M \label{up-bd}\\
           & Z_{jp} &  \in & \{0,1\} & \forall j \in M, \forall p \in P_{j} \\
           & Z_{Ai} &  \in & \{0,1\} & \forall i \in N \\
           & T_{j}  & \in & \{0,1\} & \forall  j \in M \\
           & 0 \leq & T_{A} & \leq  1 
\end{optprog} 

There is always a feasible solution if we allow the artificial columns, $T_A,Z_{A0},Z_{A1},..,Z_{A(N-1)}$
to be included, since they will eliminate any infeasibility caused by the set partitioning constraints 
(\ref{spc}) and/or  the branching constraints of (\ref{low-bd}) and (\ref{up-bd}).  

We add these $|N|+1$ variables and $|N|$ constraints to the model at the root node and define large costs for them:
\begin{align*}
 FC_A & =nREQ \cdot max_{j \in M}\{FC_j\}+1 \nonumber \\
 C_A & = \sum_{i \in N}(2 \cdot t_{facility_j,i}) \cdot OC + VF  \nonumber
\end {align*} 

The fixed cost of the artificial facility is set equal to the minimum number of required facilities 
times the maximum fixed cost of the facilities plus one to get a larger cost than selecting
real facilities. The total operating cost of each artificial 
pairing is the same with  a value equal to the sum of the operating costs of all single 
customer routes from depot $j$. We can select depot $j$ arbitrarily. The fixed cost of an artificial pairing is equal to the vehicle fixed cost. 
 



\zaupdate{\bf TAILING OFF IMPLEMENTATION}










%we created a column pool and path pool. The column pool includes
%every column created with CG algorithm. The path pool includes every path that is included in at 
%least one of the pairings in the column pool. First we branch on the location variables. When 
%there are no fractional location variables, we branch on paths. We construct a network in which 
%each path in the path pool is represented by a node. We add a source node and a sink node to 
%the network. We carry the cost and the total travel time of each path on to the incoming arcs of 
%the node representing that path. Also, each node has a demand that is equal to the total demand 
%of the path corresponding to that node. An adjacency matrix is constructed according to the customers 
%in each path. Suppose that P1=(source-1-2-sink) and P2=(source-2-3-4-sink). There is no arc 
%between P1 and P2 in the constructed network since both paths include the customer node 2. Otherwise,
%if we were to form a pairing with P1 and P2, it would not be integer feasible because of the set 
%partitioning constraint (\ref{spp_con1}). We include arcs from the source to all nodes (paths), 
%arcs from all nodes (paths) to the sink and all arcs between two path nodes that are adjacent according 
%to the adjacency matrix.
%*******************
%The ESPPRC algorithm is used over this network to find an elementary shortest path from the source to 
%the sink and new columns are created according to the branching rules. Branching rules for the paths 
%are applied with some modifications on the network. For the left branch, we set the travel time 
%between path1 and path2 to infinity. For right branch, we set the travel time on all the arcs going out 
%of path1 and all the arcs going into path2 to infinity except the arc between path1 and path2. 
%With this modification, the structure of the pricing problem does not change but the new pairings 
%satisfy the branching rules. Also, before we process the left or the right branch, we must fix 
%the values of some variables. For example, for the left branch, the pairings in which path2 follows path1
%are set to zero. For the right branch, the pairings in which there is only one of path1 or path2 or 
%both path1 and path2 but in incorrect order (not like $path1\to path2$) are fixed to zero.
%In some of our solutions a fractional pairing may include only one path. If this is the case, we may 
%apply our branching strategy to the arc between source and path1 similar to the arc between path1 and path2.
